\documentclass[10pt]{beamer}

% Tema Limpo e Profissional
\usetheme{Madrid}
\usecolortheme{whale}

% Pacotes Essenciais
\usepackage[utf8]{inputenc}
\usepackage[portuguese]{babel}
\usepackage{graphicx}
\usepackage{booktabs}
\usepackage{amsmath}
\usepackage{amssymb}
\usepackage{listings}
\usepackage{xcolor}

% Configuracoes de Codigo (Python)
\lstset{
    language=Python,
    basicstyle=\ttfamily\scriptsize,
    keywordstyle=\color{blue}\bfseries,
    stringstyle=\color{orange},
    commentstyle=\color{gray!80!black}\itshape,
    breaklines=true,
    showstringspaces=false,
    backgroundcolor=\color{gray!5},
    frame=single,
    rulecolor=\color{gray!30},
    numbers=left,
    numberstyle=\tiny\color{gray},
    stepnumber=1,
    numbersep=5pt
}

% Metadados da Apresentacao
\title[Planejamento Financeiro em Python]{Projeto Final: Planejamento Financeiro}
\subtitle{Do Orçamento Pessoal à Modelagem Computacional em Python}
\author[Prof. Reginaldo]{Prof. Reginaldo Fernandes\texorpdfstring{\\ \small{Lógica de Programação com Python}}{}}
\institute[IFCE]{Instituto Federal do Ceará -- Campus Tauá \\ Curso Técnico em Redes de Computadores}
\date{\today}

\begin{document}

% Slide 1: Titulo
\begin{frame}
    \titlepage
\end{frame}

% Slide 2: Sumario
\begin{frame}{Roteiro da Apresentação}
    \tableofcontents
\end{frame}

% ==============================================================================
% SEÇÃO 1: O PROBLEMA REAL E A REGRA 50-30-20
% ==============================================================================
\section{O Desafio do Planejamento Financeiro}

\begin{frame}{O Problema: Para Onde Vai o Dinheiro?}
    \begin{columns}
        \begin{column}{0.55\textwidth}
            \textbf{Dificuldades comuns no dia a dia:}
            \begin{itemize}
                \item Falta de visibilidade clara entre o que entra e o que sai.
                \item Gastos invisíveis (pequenas compras rotineiras que acumulam).
                \item Ausência de reserva para imprevistos e emergências.
                \item Sensação de que o dinheiro acaba antes do fim do mês.
            \end{itemize}
        \end{column}
        \begin{column}{0.45\textwidth}
            \begin{block}{O Papel da Computação}
                Softwares orçamentários transformam números brutos em \textbf{métricas visuais} e \textbf{diagnósticos automáticos}, auxiliando na tomada de decisão consciente.
            \end{block}
        \end{column}
    \end{columns}
\end{frame}

\begin{frame}{A Metodologia Orçamentária: Regra 50-30-20}
    A regra divide a receita líquida mensal em três categorias estratégicas:

    \vspace{0.3cm}
    \begin{columns}
        \begin{column}{0.33\textwidth}
            \begin{exampleblock}{50\% Essenciais}
                \textbf{Necessidades vitais:}
                \begin{itemize}
                    \item Moradia (aluguel)
                    \item Alimentação
                    \item Transporte
                    \item Contas de consumo
                    \item Saúde
                \end{itemize}
            \end{exampleblock}
        \end{column}

        \begin{column}{0.33\textwidth}
            \begin{alertblock}{30\% Estilo de Vida}
                \textbf{Desejos pessoais:}
                \begin{itemize}
                    \item Lazer e passeios
                    \item Assinaturas digitais
                    \item Restaurantes
                    \item Compras e hobbies
                \end{itemize}
            \end{alertblock}
        \end{column}

        \begin{column}{0.33\textwidth}
            \begin{block}{20\% Reserva e Metas}
                \textbf{Segurança e futuro:}
                \begin{itemize}
                    \item Reserva de emergência
                    \item Quitação de dívidas
                    \item Investimentos
                    \item Projetos de longo prazo
                \end{itemize}
            \end{block}
        \end{column}
    \end{columns}
\end{frame}

% ==============================================================================
% SEÇÃO 2: A ARQUITETURA DO SOFTWARE
% ==============================================================================
\section{Arquitetura e Fluxo do Sistema}

\begin{frame}{A Visão Geral da Solução em Python}
    O sistema foi concebido como uma aplicação CLI (Interface em Linha de Comando) dividida em 3 pilares fundamentais:

    \vspace{0.4cm}
    \begin{columns}
        \begin{column}{0.32\textwidth}
            \begin{block}{1. Entrada e Estado}
                \begin{itemize}
                    \item Previsão de Renda.
                    \item Lançamento seguro de Despesas.
                    \item Validação de tipos com \texttt{try/except}.
                \end{itemize}
            \end{block}
        \end{column}

        \begin{column}{0.32\textwidth}
            \begin{block}{2. Estrutura}
                \begin{itemize}
                    \item Dicionários (\texttt{dict}) para cada gasto.
                    \item Lista (\texttt{list}) agregadora.
                    \item Menus em ciclo (\texttt{while True}).
                \end{itemize}
            \end{block}
        \end{column}

        \begin{column}{0.32\textwidth}
            \begin{block}{3. Inteligência}
                \begin{itemize}
                    \item Somatórios por categoria.
                    \item Cálculo de saldo real.
                    \item Avaliação percentual da regra 50-30-20.
                \end{itemize}
            \end{block}
        \end{column}
    \end{columns}
\end{frame}

\begin{frame}{O Ciclo de Vida do Software (Listening Loop)}
    \begin{center}
        \textbf{Como o programa se mantém vivo no terminal?}
    \end{center}

    \begin{itemize}
        \item \textbf{Inicialização:} Configura o estado inicial na memória (\texttt{renda = 0.0}, \texttt{despesas = []}).
        \item \textbf{Laço Principal:} O comando \texttt{while True} garante que o menu não se encerre após uma ação.
        \item \textbf{Escuta:} O comando \texttt{input()} aguarda a decisão do usuário técnico.
        \item \textbf{Desvio Condicional:} \texttt{if / elif / else} encaminha para a função especialista responsável.
        \item \textbf{Término Gracioso:} A opção \texttt{"0"} executa o comando \texttt{break}, liberando o terminal.
    \end{itemize}
\end{frame}

% ==============================================================================
% SEÇÃO 3: MODELAGEM DOS DADOS
% ==============================================================================
\section{Modelagem dos Dados em Python}

\begin{frame}[fragile]{Modelagem: Dicionário como Registro Individual}
    \textbf{Analogia de Redes:} Assim como uma entrada na tabela ARP contém IP, MAC e Interface, cada despesa no software possui campos bem definidos:

    \vspace{0.2cm}
    \begin{lstlisting}
# Cada gasto do mundo real torna-se um Dicionario
despesa = {
    "descricao": "Aluguel da Residencia",
    "valor": 1200.00,
    "categoria": "Essencial"
}

# Acesso direto e legivel atraves das chaves:
print(despesa["descricao"]) # Imprime: Aluguel da Residencia
print(despesa["valor"])     # Imprime: 1200.0
    \end{lstlisting}

    \vspace{0.2cm}
    \begin{block}{Vantagem do Dicionário}
        Substitui variáveis soltas e desorganizadas por um \textbf{objeto estruturado} em que cada dado possui significado explícito por chave textual.
    \end{block}
\end{frame}

\begin{frame}[fragile]{Modelagem: Lista como Coleção de Registros}
    Para armazenar múltiplos lançamentos ao longo do mês, agrupamos os dicionários dentro de uma lista dinâmica:

    \vspace{0.2cm}
    \begin{lstlisting}
# Colecao inicial vazia
despesas = []

# Novos registros sao inseridos com .append()
despesas.append({
    "descricao": "Internet Fibra",
    "valor": 110.0,
    "categoria": "Essencial"
})
despesas.append({
    "descricao": "Cinema e Pipoca",
    "valor": 65.0,
    "categoria": "Estilo de Vida"
})

# Quantidade total armazenada:
print(len(despesas)) # Retorna: 2
    \end{lstlisting}
\end{frame}

% ==============================================================================
% SEÇÃO 4: LÓGICA DE NEGÓCIO E RELATÓRIO
% ==============================================================================
\section{Lógica de Negócio e Diagnóstico}

\begin{frame}[fragile]{O Motor de Análise: Acumuladores e Diagnóstico}
    \textbf{Como o algoritmo processa os dados na Opção 4?}

    \vspace{0.2cm}
    \begin{lstlisting}
# 1. Acumuladores condicionais com laco for:
for item in despesas:
    if item["categoria"] == "Essencial":
        total_essencial += item["valor"]
    elif item["categoria"] == "Estilo de Vida":
        total_estilo_vida += item["valor"]
    elif item["categoria"] == "Reserva":
        total_reserva += item["valor"]

# 2. Calculo do saldo liquido:
total_gasto = total_essencial + total_estilo_vida + total_reserva
saldo_restante = renda_prevista - total_gasto

# 3. Diagnostico percentual (Meta 50-30-20):
if renda_prevista > 0:
    pct_essencial = (total_essencial / renda_prevista) * 100
    \end{lstlisting}
\end{frame}

\begin{frame}{Diagnóstico Automatizado da Saúde Financeira}
    O sistema classifica a situação financeira em três estados orçamentários:

    \vspace{0.3cm}
    \begin{itemize}
        \item \textbf{Orçamento Deficitário (\texttt{saldo < 0}):}
        Os gastos superam as entradas. Alerta vermelho acionado para cortes imediatos.
        \item \textbf{Orçamento no Limite (\texttt{saldo == 0}):}
        Saldo zerado. Não há reserva para contingências.
        \item \textbf{Orçamento Superavitário (\texttt{saldo > 0}):}
        Sobra financeira. Recursos disponíveis para poupança e metas futuras.
    \end{itemize}

    \vspace{0.3cm}
    \begin{block}{Defesa Contra Divisão por Zero}
        A verificação \texttt{if renda\_prevista > 0:} impede que o Python quebre com \texttt{ZeroDivisionError} caso o usuário acesse o relatório sem registrar a receita.
    \end{block}
\end{frame}

% ==============================================================================
% SEÇÃO 5: CRONOGRAMA DE CONSTRUÇÃO
% ==============================================================================
\section{Nossa Jornada: As 5 Aulas do Projeto}

\begin{frame}{Mapa de Construção Passo a Passo}
    Construiremos o software juntos, aula a aula, em pequenas fatias funcionais:

    \vspace{0.3cm}
    \begin{enumerate}
        \item \textbf{Aula 1:} Estrutura do Programa, Loop Principal e Menu Interativo.
        \item \textbf{Aula 2:} Validação de Entrada com \texttt{try/except} e Previsão de Renda.
        \item \textbf{Aula 3:} Cadastro de Despesas com Listas e Dicionários.
        \item \textbf{Aula 4:} Listagem Formatada em Tabela e UX no Terminal.
        \item \textbf{Aula 5:} Inteligência de Negócio, Regra 50-30-20 e Entrega Final.
    \end{enumerate}

    \vspace{0.3cm}
    \begin{center}
        \textit{"Não tente escrever 200 linhas de uma vez. Programadores profissionais constroem software um bloco funcional por vez!"}
    \end{center}
\end{frame}

\begin{frame}
    \begin{center}
        \Huge{\textbf{Mãos à Obra!}}

        \vspace{0.5cm}
        \normalsize{Vamos iniciar com a \textbf{Etapa 1}: o menu e o fluxo de controle.}

        \vspace{1.0cm}
        \small{Instituto Federal do Ceará -- Campus Tauá \\ Curso Técnico em Redes de Computadores}
    \end{center}
\end{frame}

\end{document}
